\begin{problem}[第32届全国中学生物理竞赛决赛]{51}
光在物体表面反射或者被吸收时，光子将其动量传给物体，产生光的辐射压力（光压）。利用光压，可实现对某些微小粒子的精确操控（光镊）。设在 $|y|\geq L$ 的区域有一匀强激光场，沿z轴的负方向入射，其强度（单位时间内通过单位横截面积的光能）为I；在 $y\in(-L,L)$ 之间没有光场，其横截面如图所示。一个密度为 $\rho$ 的三棱柱形小物体（其横截面是底边长为2L、底角为 $\theta$ 的等腰三角形）被置于光滑的水平面xy上，其朝上的A和B两面涂有特殊反射层，能完全反射入射光。小物体初始时静止，位置如图。
\begin{subproblem}
假定光子的反射角等于入射角，且反射前后光子的频率不变。试求：
\begin{subsubproblem}
小物体在受到激光场照射后的动力学方程；
\end{subsubproblem}
\begin{subsubproblem}
小物体从初始位置向 y 轴负方向移动 3L/2 的距离所需的时间；
\end{subsubproblem}
\end{subproblem}
\begin{subproblem}
若考虑小物体运动对反射光子频率的影响，反射光子频率会有微小的变化；在小物体从初始位置向 y 轴负方向移动到距离为 3L/2 的过程中，定性比较反射光子与入射光子频率的大小及其随时间的变化。
\end{subproblem}
\end{problem}